\documentclass[12pt]{article}
\usepackage[utf8]{inputenc}
\usepackage[T1]{fontenc}
\usepackage{hyperref}
\title{Tamazight, Tidikelt Language Resources for K-12 Curriculum Development}
\author{Compiled by Claude Sonnet 5 (claude-sonnet-5) for Andr\'as Kornai}
\date{}
\begin{document}
\maketitle
\section{Introduction}
Tidikelt Tamazight (tia; Glottolog tidi1241) is a Zenati Berber language spoken around In Salah, Algeria. It is endangered. Wikipedia gives about 1,000 speakers (2011), while an older entry lists 9,000. Search turned up documentation and reference sources but no school materials.
\section{Primary Educational Resources}
None found.
\section{Linguistic and Reference Materials}
\begin{itemize}
\item Wikipedia article on the language: \url{https://en.wikipedia.org/wiki/Tidikelt_language}
\item Wikipedia article on the Tidikelt--Tuat Tamazight language: \url{https://en.wikipedia.org/wiki/Tidikelt%E2%80%93Tuat_Tamazight_language}
\item Documentation and description of Tidikelt-Tuat Tamazight (Culture in Crisis): \url{https://cultureincrisis.org/projects/documentation-and-description-of-tidikelt-tuat-tamazight-an-endangered-berber-language}
\end{itemize}
\section{Religious and Community Resources}
None found.
\section{Digital and Audio Resources}
ELDP archive record, Tidikelt-Tuat Berber: documentation of language and culture: \url{https://eldp.access.preservica.com/dk0671/}. The project, led by Samir Ben Si Said with Lameen Souag and Alexis Michaud, planned a corpus, sketch grammar and lexicon. I did not check what has been released.
\section{Curriculum Development Context}
Only older adults speak it as a first language, so any curriculum would be revitalisation work built on the documentation corpus.
\end{document}